% =====================================================================
% K3xE BKM superalgebra and Borcherds weight constants.
% =====================================================================

\providecommand{\kBKM}{\kappa_{\mathrm{BKM}}}
\providecommand{\kch}{\kappa_{\mathrm{ch}}}
\providecommand{\kcat}{\kappa_{\mathrm{cat}}}
\providecommand{\kfib}{\kappa_{\mathrm{fiber}}}
\providecommand{\Sp}{\mathrm{Sp}}
\providecommand{\OO}{\mathrm{O}}
\providecommand{\fg}{\mathfrak{g}}
\providecommand{\fn}{\mathfrak{n}}
\providecommand{\cO}{\mathcal{O}}
\providecommand{\CoHA}{\mathrm{CoHA}}
\providecommand{\mult}{\mathrm{mult}}
\providecommand{\Dfive}{D_5}
\providecommand{\Phiten}{\Phi_{10}}

\section*{K3$\times E$ BKM Superalgebra and Borcherds Weight}
\label{sec:kazhdan-K3E-BKM}

The compact Calabi--Yau threefold $X=K3\times E$ carries a
Borcherds--Kac--Moody denominator branch, a chiral specialisation
branch, a categorical Euler branch, and a fibre-lattice branch. The
four numbers
\[
  \{\kcat(X),\ \kch^{\mathrm{Heis}}(X/E),\
  \kBKM(\Delta_5),\ \kfib(K3)\}
  =
  \{0,3,5,24\}
\]
come from four constructions.

\subsection*{Primitive CHL Denominator Branch}

For the primitive CHL/BKM denominator branch attached to
$N\in\{1,2,3,4,6\}$, the Vol~III constants are
\[
\begin{array}{c|ccccc}
N & 1 & 2 & 3 & 4 & 6 \\ \hline
c_N(0) & 10 & 8 & 6 & 4 & 2 \\
\kBKM(\Phi_N)=c_N(0)/2 & 5 & 4 & 3 & 2 & 1
\end{array}
\]
Here $\Phi_N$ denotes the primitive denominator form. The scalar square
appearing in the reduced Donaldson--Thomas partition function has its
own normalisation.
At $N=1$,
\[
  c_1(0)=10,\qquad
  \kBKM(\Delta_5)=c_1(0)/2=5.
\]
Borcherds' singular theta lift gives the denominator normalisation
\[
  \operatorname{den}(\fg_{\Delta_5})=\Dfive(2Z),\qquad
  \Dfive=64^{-1}\Delta_5.
\]
The Igusa square uses
\[
  \Phiten^{\theta}=\Delta_5^2,\qquad
  \Phiten^{\mathrm{OP}}=\Dfive^2=4096^{-1}\Delta_5^2
\]
according to the theta or Oberdieck--Pandharipande scalar convention.
The weight $5$ belongs to $\Delta_5$ and $\Dfive$; the scalar square has
weight $10$.

\subsection*{K3 Elliptic Genus Normalisation}

The K3 elliptic genus is
\[
  Z_{K3}(\tau,z)=2\,\phi_{0,1}^{\mathrm{EZ}}(\tau,z),
\]
where
\[
  \phi_{0,1}^{\mathrm{EZ}}(\tau,z)
  =
  r^{-1}+10+r+O(q)
\]
in the Eichler--Zagier normalisation. The Borcherds input uses the
normalised coefficient $c_1(0)=10$; the second-quantised scalar
partition function lands on the square
$\Phiten^{\theta}=\Delta_5^2$. The fixed-locus or twining tuple
$(10,4,2,2,2)$ belongs to a separately named normalisation. The
primitive CHL row remains $(10,8,6,4,2)$.

\subsection*{Four Construction Values on $K3\times E$}

\[
\begin{array}{@{}lll@{}}
\hline
Invariant & Value & Construction \\
\hline
\kcat(X)=\chi(\cO_X) & 0 &
  \chi(\cO_{K3})\chi(\cO_E)=2\cdot 0 \\
\kch^{\mathrm{Heis}}(X/E) & 3 &
  Heisenberg--Mukai chiral specialisation over the elliptic base \\
\kBKM(\Delta_5) & 5 &
  Borcherds weight $c_1(0)/2$ of the primitive denominator \\
\kfib(K3) & 24 &
  Mukai lattice rank $\mathrm{rk}\,\widetilde H(K3,\mathbb Z)$ \\
\hline
\end{array}
\]

The compact Hodge supertrace of $X$ is
\[
  \sum_q(-1)^q h^{0,q}(K3\times E)
  =
  1-1+1-1=0,
\]
because $h^{0,\bullet}(K3\times E)=(1,1,1,1)$. Thus the compact
$\kch$ value agrees with $\kcat(X)=0$, while
$\kch^{\mathrm{Heis}}(X/E)=3$ is a fibre-relative chiral
specialisation. The fibre scalar $\chi(\cO_{K3})=2$ is a Serre-duality
witness on the K3 fibre. The total-space value is
$\kcat(K3\times E)=0$.

\subsection*{Hall--Drinfeld Branch Separation}

The $K3\times E$ BKM object is the Hall--Drinfeld double or compact Hall
branch, subject to the compactness, orientation, PBW, and finite-window
comparison hypotheses required for the source Hall construction:
\[
  \mathcal D_{\hbar}\bigl(\CoHA(K3\times E)\bigr)
  \quad\leadsto\quad
  \fg_{\Delta_5}.
\]
The positive half is the Hall algebra generated by the effective BPS
cone; the BKM denominator is recovered after the completed double and
the Borcherds product normalisation. This is the Hall--Drinfeld BKM
branch. The historical K3 Yangian belongs to the separate Mukai
self-mirror branch.

Six constructions meet at $G(K3\times E)$: the Hall algebra, stable
envelopes, Koszul duality, chiral specialisation, BKM screening, and
Borcherds lift. They construct related objects through different input
data. The single functor $\Phi$ has one output on a fixed
Calabi--Yau category.

\subsection*{Operadic Scope}

The chiral algebra $\mathbf H_{\Delta_5}$ is $E_1$-native on the
elliptic curve direction. The $E_2$ structure lives on the derived
Drinfeld centre
\[
  Z^{\mathrm{der}}_{\mathrm{ch}}(\mathbf H_{\Delta_5})
\]
and on the braided representation category produced by the
Maulik--Okounkov stable-envelope $R$-matrix. At the hyper-Kahler limit
the $R$-matrix is symmetric; the non-trivial braided structure appears
after the $\Omega$-background deformation.

\subsection*{Primary Sources}

\begin{itemize}[leftmargin=1.5em]
\item Borcherds 1998, \emph{Inventiones}~132, Theorem~13.3: singular
  theta lift weight $\mathrm{wt}(\Phi)=c_0(0)/2$.
\item Gritsenko--Nikulin 1995/1998: denominator
  $\operatorname{den}(\fg_{\Delta_5})=\Dfive(2Z)$ with
  $\Dfive=64^{-1}\Delta_5$.
\item Eichler--Zagier 1985: $\phi_{0,1}^{\mathrm{EZ}}$ and its
  coefficient $c_1(0)=10$.
\item Oberdieck--Pandharipande 2016 and Oberdieck--Pixton 2018:
  $K3\times E$ reduced DT partition functions and the Igusa square.
\item Maulik--Okounkov 2019 and Kontsevich--Soibelman 2011:
  stable-envelope $R$-matrices and CoHA $E_1$ structure.
\end{itemize}
